\documentclass[12pt]{article}
\usepackage[utf8]{inputenc}
\usepackage[T1]{fontenc}
\usepackage{hyperref}
\title{Papuan Malay Language Resources for K-12 Curriculum Development}
\author{Compiled by Claude Sonnet 5.5 (claude-sonnet-5-5) for Andr\'as Kornai}
\date{}
\begin{document}
\maketitle
\section{Introduction}
Papuan Malay (ISO 639-3 pmy) is a Malay-based lingua franca of West Papua, Indonesia. It is not officially recognised for formal government or educational use. Searches found scholarly description but no school materials.
\section{Primary Educational Resources}
No primers, textbooks or school programmes in Papuan Malay were found.
\section{Linguistic and Reference Materials}
\begin{itemize}
\item A grammar of Papuan Malay (Angela Kluge), Language Science Press, open access: \url{https://langsci-press.org/catalog/book/78}
\item The Leiden University open-access repository record for the same work (2014 dissertation): \url{https://openaccess.leidenuniv.nl/handle/1887/25849}
\item A 2026 conference paper (Schapper) describing Papuan Malay as among the best described eastern Indonesian Malay varieties: \url{https://events.tuni.fi/app/uploads/sites/1396/2026/06/schapper.pdf}
\end{itemize}
No dictionary was found in searches.
\section{Religious and Community Resources}
None found in searches.
\section{Digital and Audio Resources}
None confirmed. A Papuan Malay entry in the CELD Papua collection (Himmelmann and Riesberg) is cited in the paper above, but no direct URL was seen.
\section{Curriculum Development Context}
Since the variety is a second-language lingua franca with no official educational status, any curriculum would start from the Kluge grammar and would need new materials.
\end{document}